% =============================================================================
\chapter{Auditoria de qualidade e ETL}\label{cap:etl}
% =============================================================================

Este capítulo documenta o \textit{Checkpoint}~2. A ordem das seções segue a do
notebook \code{A3_CP2_auditoria_etl_eda.ipynb}: extração, auditoria do dado
bruto (sem alterar nada, apenas medir), transformação, validação cruzada com a
fonte primária e carga dos produtos analíticos. Cada problema encontrado na
auditoria foi registrado em um log estruturado (\code{AUDIT_LOG}), e cada etapa
da transformação, em uma tabela de linhagem (\code{LINEAGE}).

\section{Extração com verificação de integridade}\label{sec:extracao}

O arquivo é baixado de um endereço que fixa o \textit{commit} do repositório no
Hugging Face e, em seguida, conferido por SHA-256 (\autoref{qua:ficha}). Se um
único \textit{byte} estiver diferente, por download corrompido ou alteração do
arquivo pelo fornecedor, a execução é interrompida e o arquivo inválido é
apagado, em vez de seguir com dados errados. Esse é o primeiro filtro contra o
\textit{garbage in}. O download grava primeiro um arquivo temporário
(\code{.part}) e só o renomeia ao final, para nunca deixar um arquivo pela
metade com o nome definitivo. Na execução registrada, o arquivo de 103,4~MB foi
baixado em 16,9~s e carregado em 17,3~s, ocupando 1,80~GB de memória.

Antes de carregar os dados, os metadados do Parquet já revelaram dois sinais de
alerta: \code{trade_id} armazenado como texto, embora seja um número
sequencial, e os dois carimbos de tempo armazenados como inteiros, não como
data-hora.

\section{Auditoria do dado bruto}\label{sec:auditoria}

\subsection{Completude}\label{sec:completude}

A completude foi verificada em três níveis: nulos explícitos, nulos disfarçados
(textos vazios, \code{"null"}, \code{"N/A"}, zeros ou negativos onde não fazem
sentido, identificadores não numéricos) e registros faltantes. Em um fluxo de
negócios, um nulo também pode ser um negócio que deveria existir e não está no
arquivo; como o \code{trade_id} da Binance é sequencial, buracos na sequência
revelariam perdas.

O resultado foi zero nulos explícitos, zero nulos disfarçados e
\textbf{nenhum negócio faltante}: os 27.925.338 identificadores esperados
($\max - \min + 1$) estão todos presentes, e os 10.080 minutos da semana têm
pelo menos um negócio, com intervalo máximo de 8,2~s entre dois negócios
consecutivos. O único problema de completude está nos \textbf{metadados}: o nome
do arquivo indica o período até 26/02, mas o último negócio é de 25/02 às
23:59:59~UTC; a data final do nome é exclusiva.

\subsection{Unicidade e duplicatas aparentes}\label{sec:unicidade}

Foram aplicados três testes, do mais estrito ao mais frouxo: linha inteira
repetida (zero ocorrências), chave \code{trade_id} repetida (zero ocorrências)
e \textbf{duplicata aparente}, isto é, linhas iguais em tudo exceto o
\code{trade_id}. O terceiro teste marcou \textbf{11.631.537 linhas, 41,7\,\% da
base}.

A inspeção dos casos mostrou que essas linhas são o funcionamento normal do
motor de negociação (\autoref{sec:microestrutura}): uma ordem agressora casada
contra várias ordens passivas de mesmo preço e tamanho gera vários negócios no
mesmo milissegundo, com identificadores consecutivos. A \autoref{tab:dropdup}
quantifica o que aconteceria com uma remoção ingênua.

\begin{table}[htbp]
\caption{\label{tab:dropdup}Impacto de uma remoção ingênua de duplicatas aparentes}
\begin{corpotabela}
\begin{tabularx}{\textwidth}{L{3.6cm}RRRR}
\toprule
Operação & Negócios apagados & \% dos negócios & Volume apagado (\usd~mi) & \% do volume \\
\midrule
\code{drop_duplicates()} ingênuo & 11.631.537 & 41,7 & 410,5 & 2,2 \\
\bottomrule
\end{tabularx}
\end{corpotabela}
\fontedados
\end{table}

A maioria das duplicatas aparentes são negócios minúsculos: a quantidade de
0,00006~BTC (cerca de \usd~5,7, próximo do valor mínimo de ordem da Binance)
aparece 10.156.068 vezes entre elas. \textbf{Decisão: manter.} Apagá-las
distorceria qualquer métrica de contagem (P1 e P2) e o fluxo de ordens (P3). Em
vez disso, o ETL cria uma tabela de ordens agressoras que agrupa as execuções
parciais (\autoref{sec:ordens}), e a validação cruzada
(\autoref{sec:validacaocruzada}) prova que essas linhas existem na fonte
oficial.

\subsection{Validade}\label{sec:validade}

Os domínios foram verificados por regras explícitas: preços positivos e
múltiplos do \textit{tick} de 0,01~USDT, quantidades positivas e múltiplas do
passo de 0,00001~BTC, lados pertencentes a $\{-1, +1\}$ e
$\{\text{buy}, \text{sell}\}$, carimbos dentro do período declarado. Todas
passaram, com duas exceções de projeto: \code{exchange} e \code{symbol} têm
cardinalidade~1, ou seja, não informam nada, e \code{trade_id} ocupa 503~MB como
texto contra 223~MB como inteiro.

A regra mais importante é a do carimbo de ingestão. Por definição,
\code{ingest_ts} é o instante em que o fornecedor gravou o negócio; portanto
deve ser posterior ao negócio e anterior à publicação do arquivo
(25/10/2026). A \autoref{tab:ingest} mostra que a coluna tem apenas sete valores,
um por dia de negociação, todos entre 2032 e 2037.

\begin{table}[htbp]
\caption{\label{tab:ingest}Valores distintos de \texttt{ingest\_ts} na base bruta}
\begin{corpotabela}
\begin{tabular}{llrr}
\toprule
\code{ingest_ts} convertido (UTC) & Dia dos negócios & Negócios & Atraso (anos) \\
\midrule
05/04/2032 16:40:00 & 19/02/2025 & 2.784.225 & 7,1 \\
15/12/2032 23:10:00 & 20/02/2025 & 2.386.973 & 7,8 \\
21/08/2033 06:20:00 & 21/02/2025 & 4.528.525 & 8,5 \\
23/07/2034 15:23:20 & 22/02/2025 & 1.506.418 & 9,4 \\
02/01/2035 09:30:00 & 23/02/2025 & 1.457.368 & 9,9 \\
07/07/2035 03:00:00 & 24/02/2025 & 4.972.684 & 10,4 \\
09/01/2037 23:20:00 & 25/02/2025 & 10.289.145 & 11,9 \\
\bottomrule
\end{tabular}
\end{corpotabela}
\fontedados
\end{table}

Nenhum registro passa na regra: \code{ingest_ts} está \textbf{100\,\% inválido}.
É um campo corrompido no pipeline do fornecedor, provavelmente por unidade ou
cálculo errado. Usado sem auditoria, por exemplo para medir a latência de
coleta, produziria conclusões absurdas; é o caso clássico de \textit{Garbage
In, Garbage Out}.

\subsection{Consistência}\label{sec:consistencia}

As colunas \code{raw_side} e \code{aggressor_side} concordam em 100\,\% das
linhas, o que caracteriza redundância. A ordem dos identificadores e a ordem
temporal também são consistentes: \code{trade_id} é estritamente crescente e
\code{exchange_ts} é não decrescente, ou seja, nenhum negócio volta no tempo.

\subsection{Exatidão e valores extremos}\label{sec:outliers}

A quantidade negociada tem assimetria de 92,8 e curtose de 49.183. Nessas
condições, a regra clássica da amplitude interquartil ($Q_3 + 1,5 \cdot
\mathrm{IQR}$) marcaria como \textit{outlier} qualquer negócio acima de
0,00064~BTC (cerca de \usd~60), isto é, \textbf{20,0\,\% de todos os negócios}.
A regra pressupõe uma distribuição aproximadamente simétrica, o que não vale
aqui. A \autoref{tab:top5} mostra os cinco maiores negócios: são reais e são
justamente os que mais importam para a pergunta P2.

\begin{table}[htbp]
\caption{\label{tab:top5}Os cinco maiores negócios da semana, por valor financeiro}
\begin{corpotabela}
\begin{tabular}{rlrrlr}
\toprule
\code{trade_id} & Instante (UTC) & Preço (USDT) & Quantidade (BTC) & Lado & Valor (\usd) \\
\midrule
4.581.643.844 & 22/02/2025 19:05:43,614 & 96.419,17 & 56,34875 & venda & 5.433.099,71 \\
4.590.361.579 & 25/02/2025 06:17:45,712 & 91.675,99 & 37,32239 & venda & 3.421.567,05 \\
4.572.878.642 & 19/02/2025 19:00:26,039 & 95.625,00 & 24,82016 & venda & 2.373.427,80 \\
4.581.688.503 & 22/02/2025 19:57:35,018 & 96.644,99 & 21,87568 & venda & 2.114.174,87 \\
4.598.452.092 & 25/02/2025 23:05:54,389 & 89.050,32 & 21,50070 & compra & 1.914.644,22 \\
\bottomrule
\end{tabular}
\end{corpotabela}
\fontedados
\end{table}

O preço não apresenta saltos impossíveis (\textit{bad ticks}): o maior salto
entre dois negócios consecutivos é de 0,187\,\%, durante a queda de 25/02, e não
há nenhum salto acima de 1\,\%. \textbf{Decisão: não remover valores
extremos}; em seu lugar, usar escala logarítmica nos gráficos, medianas quando a
cauda distorce a média e faixas de tamanho (\code{size_bucket}).

\subsection{Síntese da auditoria}\label{sec:sinteseauditoria}

O \autoref{qua:auditoria} consolida os doze registros do log de auditoria, com a
dimensão de qualidade, a evidência, a severidade e o tratamento definido.

\begin{quadro}[htbp]
\caption{\label{qua:auditoria}Achados da auditoria de qualidade do dado bruto}
\begin{corpotabela}
\begin{tabularx}{\textwidth}{|C{0.5cm}|L{1.9cm}|L{2.4cm}|Y|L{1.6cm}|L{3.6cm}|}
\hline
\textbf{N.} & \textbf{Dimensão} & \textbf{Coluna} & \textbf{Achado e evidência} & \textbf{Severidade} & \textbf{Tratamento} \\ \hline
1 & Completude & todas & Nenhum nulo explícito ou disfarçado (0 NaN; 0 vazios, ``null'' ou $\leq 0$) & -- & Nenhum \\ \hline
2 & Completude & \code{trade_id} & Sequência completa: 27.925.338 IDs esperados e presentes & -- & Nenhum \\ \hline
3 & Consistência & metadados & Período declarado (até 26/02) diferente do observado (até 25/02 23:59:59) & baixa & Documentar o período real \\ \hline
4 & Unicidade & todas / \code{trade_id} & Sem duplicatas reais (0 linhas e 0 IDs repetidos) & -- & Deduplicação por chave como proteção idempotente \\ \hline
5 & Unicidade & todas exceto \code{trade_id} & Duplicatas aparentes (execuções parciais): 11.631.537 linhas (41,7\,\%), 2,2\,\% do volume & alta (armadilha) & Manter; agregar em ordens agressoras \\ \hline
6 & Validade & \code{ingest_ts} & Carimbos no futuro (2032 a 2037); 0\,\% passam na regra & crítica & Converter em nulo e descartar a coluna \\ \hline
7 & Validade & \code{trade_id} & Número guardado como texto (503 MB contra 223 MB) & média & Converter para inteiro de 64 bits \\ \hline
8 & Validade & \code{exchange_ts} & Data-hora guardada como inteiro em ms & média & Converter para data-hora UTC (\code{ts}) \\ \hline
9 & Validade & \code{exchange}, \code{symbol} & Colunas constantes (cardinalidade 1) & baixa & Mover para metadados e remover \\ \hline
10 & Consistência & \code{raw_side} $\times$ \code{aggressor_side} & Colunas redundantes (concordância de 100\,\%) & baixa & Manter uma só coluna categórica, \code{side} \\ \hline
11 & Exatidão & \code{quantity}, valor & Cauda extremamente pesada: IQR marcaria 20\,\%; maior negócio de \usd~5,4 mi & informativa & Não remover; usar escala log, medianas e faixas \\ \hline
12 & Exatidão & \code{price} & Sem \textit{bad ticks}: maior salto de 0,187\,\% & -- & Nenhum \\ \hline
\end{tabularx}
\end{corpotabela}
\fonteautores[Elaborado]
\end{quadro}

Dois achados adicionais surgiram nas etapas seguintes e foram registrados no
mesmo log: a divergência de unidade de tempo entre os candles oficiais e os
negócios (\autoref{sec:validacaocruzada}) e os nulos estruturais ao reamostrar
em barras de 1 segundo (\autoref{sec:gaps}).

\section{Transformação}\label{sec:transformacao}

Para não estourar a memória do Google Colab gratuito, a transformação é feita
sobre o próprio \textit{DataFrame} carregado. O dado bruto continua intacto em
disco (\code{data/raw/}, com resumo conferido) e pode ser relido a qualquer
momento.

\subsection{Adequação de tipos}\label{sec:tipos}

\begin{quadro}[htbp]
\caption{\label{qua:tipos}Conversões de tipo aplicadas}
\begin{corpotabela}
\begin{tabularx}{\textwidth}{|L{3.4cm}|L{2.2cm}|L{3.0cm}|Y|}
\hline
\textbf{Coluna} & \textbf{De} & \textbf{Para} & \textbf{Justificativa} \\ \hline
\code{trade_id} & texto & inteiro de 64 bits & Ocupa menos memória, ordena corretamente (como texto, ``10'' $<$ ``9'') e permite checar a sequência \\ \hline
\code{exchange_ts} $\rightarrow$ \code{ts} & inteiro (ms) & data-hora UTC & Permite reamostrar, extrair hora e dia e é legível \\ \hline
\code{aggressor_side} $\rightarrow$ \code{side} & inteiro ($\pm 1$) & categórica (\code{buy}/\code{sell}) & Rótulo explícito; a categoria ocupa 1 \textit{byte} por linha \\ \hline
\end{tabularx}
\end{corpotabela}
\fonteautores[Elaborado]
\end{quadro}

\subsection{Valores inválidos e estratégia para nulos}\label{sec:nulos}

O procedimento padrão para um valor que viola uma regra de validade tem dois
passos: convertê-lo em nulo explícito, para que nenhuma conta o use por engano,
e então decidir o que fazer com o nulo. Após a invalidação, \code{ingest_ts}
ficou com 27.925.338 nulos (100\,\%). O \autoref{qua:estrategias} compara as
estratégias clássicas para esse caso.

\begin{quadro}[htbp]
\caption{\label{qua:estrategias}Estratégias avaliadas para a coluna \texttt{ingest\_ts}}
\begin{corpotabela}
\begin{tabularx}{\textwidth}{|L{4.6cm}|Y|C{1.8cm}|}
\hline
\textbf{Estratégia} & \textbf{Aplicação neste caso} & \textbf{Veredito} \\ \hline
Remover linhas com nulo & Apagaria a base inteira & rejeitada \\ \hline
Imputar (média, mediana, último valor, cópia de \code{ts}) & Não há nenhum valor válido de referência; copiar \code{ts} fabricaria um dado que não existe & rejeitada \\ \hline
Criar indicador e manter & O indicador seria constante (100\,\% inválido) e não informaria nada & rejeitada \\ \hline
Remover a coluna & Nenhuma pergunta de negócio depende do instante de ingestão & adotada \\ \hline
\end{tabularx}
\end{corpotabela}
\fonteautores[Elaborado]
\end{quadro}

\subsection{Colunas constantes e redundantes}\label{sec:constantes}

As colunas \code{exchange} e \code{symbol} viram metadados do
\textit{DataFrame} (\code{df.attrs}), junto com o endereço da fonte e o SHA-256,
em vez de 27,9 milhões de cópias da mesma palavra. As colunas \code{raw_side} e
\code{aggressor_side} são substituídas pela coluna \code{side}.

\subsection{Deduplicação idempotente}\label{sec:dedup}

As duplicatas são removidas apenas pela chave \code{trade_id}. Hoje a operação
remove zero linhas, mas torna o pipeline \textbf{idempotente}: se o arquivo for
baixado duas vezes e concatenado por engano, por exemplo, o problema é corrigido
automaticamente. As duplicatas aparentes são preservadas.

\subsection{Novas variáveis}\label{sec:novasvariaveis}

São criadas duas variáveis no nível do negócio. O valor financeiro,
$\mathit{notional}_i = p_i \times q_i$, é a medida correta de tamanho, porque
1~BTC a \usd~86 mil é menos dinheiro que 1~BTC a \usd~99 mil. A faixa de tamanho
(\code{size_bucket}) é uma categórica ordenada com cinco níveis, intervalos
fechados à esquerda, usada nos agrupamentos da pergunta P2.

\subsection{Contrato de qualidade}\label{sec:contrato}

Um contrato de dados é um conjunto de regras que a base tratada tem de cumprir.
A função \code{validate_trades()} verifica as dez regras do
\autoref{qua:contrato} e interrompe a execução (\code{assert}) se qualquer uma
falhar, de modo que nenhuma análise rode sobre dado ruim. Na base tratada,
todas passaram.

\begin{quadro}[htbp]
\caption{\label{qua:contrato}Regras do contrato de qualidade da base tratada}
\begin{corpotabela}
\begin{tabularx}{\textwidth}{|C{0.6cm}|Y|C{2.0cm}|}
\hline
\textbf{N.} & \textbf{Regra} & \textbf{Resultado} \\ \hline
1 & Nenhum valor nulo em nenhuma coluna & aprovada \\ \hline
2 & \code{trade_id} único & aprovada \\ \hline
3 & \code{trade_id} contíguo, sem buracos (diferença sempre igual a 1) & aprovada \\ \hline
4 & \code{ts} em ordem cronológica & aprovada \\ \hline
5 & \code{ts} dentro de [19/02/2025, 26/02/2025) UTC & aprovada \\ \hline
6 & \code{price} $> 0$ & aprovada \\ \hline
7 & \code{quantity} $> 0$ & aprovada \\ \hline
8 & \code{side} $\in \{\text{buy}, \text{sell}\}$ & aprovada \\ \hline
9 & \code{notional_usd} $= $ \code{price} $\times$ \code{quantity} & aprovada \\ \hline
10 & Tipos esperados: inteiro de 64 bits, data-hora e categórica & aprovada \\ \hline
\end{tabularx}
\end{corpotabela}
\fonteautores[Elaborado]
\end{quadro}

\subsection{Linhagem: antes e depois}\label{sec:linhagem}

A \autoref{tab:linhagem} reproduz a tabela de linhagem gerada pelo notebook. O
ETL \textbf{não perdeu nenhum negócio}: as 27.925.338 linhas foram preservadas.
Foram removidas cinco colunas sem informação útil, ajustados três tipos e
criadas duas variáveis de análise. A memória caiu cerca de 35\,\%, de 1,80~GB
para 1,17~GB, mesmo com as duas colunas novas.

\begin{table}[htbp]
\caption{\label{tab:linhagem}Linhagem da transformação da base de negócios}
\begin{corpotabela}
\begin{tabularx}{\textwidth}{L{5.2cm}P{1.3cm}P{1.7cm}P{1.7cm}Y}
\toprule
Etapa & Colunas & Memória (MB) & Nulos & Observação \\
\midrule
0. Bruto & 9 & 1.801,7 & 0 & -- \\
1. Tipos ajustados & 10 & 1.550,4 & 0 & \code{trade_id}, \code{ts} e \code{side} \\
2a. \code{ingest_ts} inválido $\rightarrow$ nulo & 10 & 1.550,4 & 27.925.338 & 100\,\% nulos \\
2b. Remoção de \code{ingest_ts} & 9 & 1.327,0 & 0 & coluna 100\,\% inválida \\
3. Remoção de constantes e redundantes & 5 & 921,5 & 0 & metadados em \code{attrs} \\
4. Duplicatas por chave & 5 & 921,5 & 0 & 0 removidas; aparentes mantidas \\
5. Novas variáveis & 7 & 1.172,9 & 0 & \code{notional_usd}, \code{size_bucket} \\
\bottomrule
\end{tabularx}
\end{corpotabela}
\fonte{Elaborada pelos autores (2026), com dados de \citeonline{glitchlist2026}. Todas as etapas mantêm 27.925.338 linhas.}
\end{table}

\section{Validação cruzada com a fonte oficial}\label{sec:validacaocruzada}

O fornecedor da base é um terceiro. Para provar que a base é confiável e
completa, o ETL reconstrói candles de 1 minuto a partir dos negócios tratados e
os compara com os candles oficiais publicados pela própria Binance
\cite{binance2025}. Os sete arquivos diários são baixados com o SHA-256 que a
Binance publica ao lado de cada arquivo.

Um detalhe técnico precisou ser tratado antes da junção: os arquivos oficiais de
2025 usam carimbos em \textbf{microssegundos}, enquanto os negócios usam
\textbf{milissegundos}. Sem normalização, nenhum minuto casaria. A função de
conversão detecta a unidade pela ordem de grandeza do valor (segundos,
milissegundos, microssegundos ou nanossegundos). A \autoref{tab:validacao}
mostra o resultado.

\begin{table}[htbp]
\caption{\label{tab:validacao}Concordância entre os candles reconstruídos e os candles oficiais da Binance}
\begin{corpotabela}
\begin{tabular}{lr}
\toprule
Critério de comparação & Minutos concordantes (\%) \\
\midrule
Minuto presente nas duas fontes & 100,00 \\
Preço de abertura idêntico & 100,00 \\
Preço máximo idêntico & 100,00 \\
Preço mínimo idêntico & 100,00 \\
Preço de fechamento idêntico & 100,00 \\
Número de negócios idêntico & 100,00 \\
Volume em BTC igual (tolerância relativa de $10^{-9}$) & 100,00 \\
Volume em dólares igual (tolerância relativa de $10^{-9}$) & 100,00 \\
Volume comprador igual (tolerância relativa de $10^{-9}$) & 100,00 \\
\bottomrule
\end{tabular}
\end{corpotabela}
\fonte{Elaborada pelos autores (2026), com dados de \citeonline{glitchlist2026} e \citeonline{binance2025}. Base de comparação: 10.080 minutos.}
\end{table}

Todos os 10.080 minutos coincidem, e o total de negócios é igual nas duas
fontes (27.925.338). As diferenças de volume ficam na ordem de $10^{-13}$~BTC,
arredondamento de ponto flutuante. Três conclusões decorrem disso: a base do
fornecedor é fiel à Binance; nenhum negócio foi perdido, nem pelo fornecedor nem
pelo ETL; e as duplicatas aparentes existem na fonte oficial, de modo que
apagá-las teria sido um erro.

\section{Carga e agregações}\label{sec:carga}

\subsection{Barras de 1 minuto e de 1 segundo}\label{sec:barras}

A função \code{build_bars()} agrupa os negócios pelo início do intervalo
(\code{ts.dt.floor(freq)}) e calcula abertura, máxima, mínima, fechamento,
volume em BTC e em dólares, número de negócios e volume comprador. O resultado é
reindexado sobre uma \textbf{grade temporal completa}, de modo que intervalos sem
negócio aparecem explicitamente como nulos. A função \code{add_bar_features()}
acrescenta as variáveis do \autoref{qua:derivadas}. As barras de 1 minuto não
têm nenhuma lacuna; as de 1 segundo têm, e exigem tratamento.

\subsection{Nulos estruturais nas barras de 1 segundo}\label{sec:gaps}

Dos 604.800 segundos da semana, \textbf{40.202 (6,6\,\%) não têm nenhum
negócio}. Esses nulos são estruturais, não erro de coleta: o nulo significa
``não houve negócio''. As lacunas são curtas (34.732 sequências, com mediana de
1~s e máximo de 7~s) e concentram-se na madrugada: 11,5\,\% dos segundos das
06h~UTC estão vazios, contra 2,5\,\% às 15h~UTC. O \autoref{qua:gaps} mostra o
tratamento de cada coluna.

\begin{quadro}[htbp]
\caption{\label{qua:gaps}Tratamento dos nulos estruturais nas barras de 1 segundo}
\begin{corpotabela}
\begin{tabularx}{\textwidth}{|L{4.0cm}|L{4.2cm}|Y|}
\hline
\textbf{Coluna} & \textbf{Tratamento} & \textbf{Justificativa} \\ \hline
\code{close} & último preço conhecido (\code{ffill}) & O preço vigente é o do último negócio, convenção usual em finanças \\ \hline
\code{open}, \code{high}, \code{low} & iguais ao \code{close} preenchido & Barra ``plana'': o preço não se moveu \\ \hline
\code{volume_btc}, \code{volume_usd}, \code{n_trades}, \code{buy_volume_btc} & zero & Sem negócio significa volume zero, não valor desconhecido \\ \hline
\code{vwap}, \code{ofi}, \code{avg_trade_usd} & \code{close}, 0 e 0 & Valores neutros para razões indefinidas (0/0) \\ \hline
\code{has_trades} (nova) & indicador booleano & Preserva a informação de que o valor foi preenchido \\ \hline
\end{tabularx}
\end{corpotabela}
\fonteautores[Elaborado]
\end{quadro}

Duas alternativas foram rejeitadas. A \textbf{interpolação linear} usaria o
preço futuro para preencher o passado (viés de antecipação, ou
\textit{look-ahead bias}) e contaminaria qualquer modelo preditivo. A
\textbf{remoção dos segundos vazios} quebraria a grade regular de tempo exigida
por médias móveis e defasagens. O \autoref{qua:fillgaps} mostra o trecho de
código que implementa o preenchimento causal; o teste automatizado
correspondente está descrito no \autoref{cap:reprodutibilidade}.

\begin{quadro}[htbp]
\caption{\label{qua:fillgaps}Trecho de \texttt{src/etl.py}: preenchimento de lacunas sem viés de antecipação}
\begin{minipage}{\textwidth}
\SingleSpacing
\begin{lstlisting}[language=Python]
def fill_bar_gaps(bars: pd.DataFrame) -> pd.DataFrame:
    filled = bars.copy()
    filled["has_trades"] = filled["n_trades"].notna()
    for col in ("volume_btc", "volume_usd", "n_trades", "buy_volume_btc"):
        filled[col] = filled[col].fillna(0)
    filled["close"] = filled["close"].ffill()       # só olha para trás
    for col in ("open", "high", "low"):
        filled[col] = filled[col].fillna(filled["close"])
    filled["n_trades"] = filled["n_trades"].astype("int64")
    return filled
\end{lstlisting}
\end{minipage}
\fonte{Repositório do projeto (2026).}
\end{quadro}

\subsection{Ordens agressoras}\label{sec:ordens}

Em vez de apagar as execuções parciais, o ETL as agrupa. Negócios consecutivos
com o mesmo milissegundo e o mesmo lado vêm, quase certamente, da mesma ordem
agressora; a Binance usa critério semelhante em seu arquivo oficial de negócios
agregados. A função \code{build_taker_orders()} marca o início de uma nova ordem
sempre que o carimbo ou o lado mudam em relação à linha anterior e soma
quantidade e valor de cada grupo.

O resultado são \textbf{4.555.252 ordens agressoras}, com média de 6,1
execuções por ordem; 65,4\,\% das ordens têm uma única execução, e a maior
reuniu 2.024 execuções e \usd~7,2 milhões. Trata-se de uma aproximação: duas
ordens diferentes, de participantes distintos, no mesmo milissegundo e no mesmo
lado seriam fundidas. Para medir o perfil de tamanho das ordens, o erro é
desprezível.

\subsection{Produtos gravados}\label{sec:produtos}

Os quatro produtos analíticos são gravados em Parquet com compressão ZSTD em
\code{data/processed/} (\autoref{tab:produtos}). O dicionário de cada arquivo
está no \autoref{apendice:dicionario}. O painel ainda grava um quinto arquivo,
\code{trade_size_hourly.parquet}, como \textit{cache} da agregação por faixa de
tamanho.

\begin{table}[htbp]
\caption{\label{tab:produtos}Arquivos gerados pela etapa de carga}
\begin{corpotabela}
\begin{tabular}{llrrr}
\toprule
Arquivo & Granularidade & Linhas & Colunas & Tamanho (MB) \\
\midrule
\code{trades_clean.parquet} & 1 linha por negócio & 27.925.338 & 7 & 244,8 \\
\code{bars_1min.parquet} & 1 linha por minuto & 10.080 & 15 & 1,2 \\
\code{bars_1s.parquet} & 1 linha por segundo & 604.800 & 16 & 35,7 \\
\code{taker_orders.parquet} & 1 linha por ordem agressora & 4.555.252 & 10 & 89,2 \\
\bottomrule
\end{tabular}
\end{corpotabela}
\fontedados
\end{table}
